\chapter{Labia Changes (Atrophy)}

\section*{Definition}
Structural changes to the labia majora and minora including shrinkage, loss of elasticity, colour changes, and reduced tissue volume, resulting from hormonal insufficiency.

\section*{Pathophysiology}
Labial tissues contain oestrogen and androgen receptors that maintain collagen density, vascularity, and elastic fibres. Declining oestrogen during menopause causes dermal thinning, reduced collagen synthesis, loss of subcutaneous fat, and decreased blood flow. The labia minora may become pale, thin, and gradually resorbed; the labia majora lose subcutaneous fat volume. These changes contribute to vulvovaginal atrophy (part of genitourinary syndrome of menopause) and can affect urinary and sexual function.

\section*{Causes}
\begin{itemize}
  \item Genitourinary syndrome of menopause (oestrogen deficiency)
  \item Natural ageing (collagen loss independent of hormones)
  \item Lichen sclerosus (autoimmune, causes scarring and fusion)
  \item Lichen planus (inflammatory dermatosis)
  \item Crohn disease (vulvar involvement)
  \item Radiation therapy to the pelvic region
  \item Surgical menopause (bilateral oophorectomy)
  \item Chemotherapy-induced ovarian failure
  \item Autoimmune oophoritis
  \item Smoking (accelerated oestrogen metabolism)
\end{itemize}

\section*{When to See a Doctor}
See your doctor if:
\begin{itemize}
  \item Severe or worsening symptoms
  \item Labial changes with itching, burning, pain, skin thickening or whitening, fissures, or difficulty with intercourse
  \item Accompanied by other concerning signs
\end{itemize}

\section*{Related Symptoms}
\begin{itemize}
  \item Vaginal dryness
  \item Vaginal atrophy
  \item Painful sex
  \item Clitoral pain
  \item Urinary symptoms
\end{itemize}

\textbf{Important:} If this symptom persists, your doctor will check for serious underlying causes.

\section*{Self-Care / Home Management}
\begin{itemize}
  \item Rest and avoid activities that make it worse.
  \item Maintain adequate hydration and a balanced diet.
  \item Over-the-counter remedies may provide symptomatic relief (consult a pharmacist or doctor).
  \item Monitor symptoms and seek professional advice if they persist or worsen.
  \item Avoid self-medicating with prescription medications without medical guidance.
\end{itemize}

\section*{How Your Doctor Will Evaluate You}
\begin{itemize}
  \item A thorough history and physical examination are the first steps in evaluation.
  \item Laboratory tests (blood work, urinalysis) may be ordered based on clinical suspicion.
  \item Imaging studies (X-ray, ultrasound, CT, MRI) may be indicated when structural pathology is suspected.
  \item Specialist referral is arranged when the diagnosis remains uncertain or requires specific expertise.
\end{itemize}

\section*{Treatment Options}
\begin{itemize}
  \item Treatment targets the underlying cause identified through diagnostic evaluation.
  \item Supportive care and symptom management are provided while awaiting definitive therapy.
  \item Medications, lifestyle modifications, and physical therapies may be prescribed as appropriate.
  \item Follow-up ensures treatment effectiveness and allows timely adjustments to the care plan.
\end{itemize}

\clearpage